% !TEX program = pdflatex
% ------------------------------------------------------------
% MATH 431 Homework Template (plug-and-chug friendly)
%
% Goals:
% - Same clean header style as your MATH 521 template.
% - Easy “Problem -> Solution” workflow.
% - Lightweight (no tcolorbox); uses mdframed for optional boxes.
% - Good for probability homework: sample space, probability measure, events.
%
% Quick use:
%   1) Edit the Metadata block.
%   2) For each exercise, use \problem{<section>.<number>} and paste the prompt.
%   3) Write your solution in the solution environment.
% ------------------------------------------------------------

\documentclass[12pt]{article}

% =========================
% 0) Metadata (EDIT ME)
% =========================
\newcommand{\CourseCode}{MATH 431}
\newcommand{\CourseTitle}{Introduction to Probability}
\newcommand{\CourseSection}{LEC 002} % change to 001 or 003 if needed
\newcommand{\Semester}{Spring 2026}
\newcommand{\Instructor}{Sarah Tammen} % LEC 001: Ander Aguirre | LEC 002: Sarah Tammen | LEC 003: Nicholas Derr
\newcommand{\MeetingInfo}{MWF 9:55--10:45, 6102 Sewell Social Sciences} % LEC 001 default; update if in 002/003
\newcommand{\HomeworkNumber}{1}
\newcommand{\YourName}{Alejandro De La Torre}
\newcommand{\DueDate}{February 20, 2026} % or e.g. February 2, 2026

% =========================
% 1) Core math tools
% =========================
\usepackage{amsmath, amssymb, amsthm}

% =========================
% 2) Lists and spacing
% =========================
\usepackage{enumitem}
\setlist[enumerate,1]{label=\textbf{(\alph*)}, leftmargin=2em, itemsep=0.25em}
\setlist[itemize]{leftmargin=1.5em, itemsep=0.25em}
\setlength{\parskip}{0.35em}
\setlength{\parindent}{0pt}

% =========================
% 3) Page geometry + header/footer
% =========================
\usepackage{geometry}
\geometry{margin=1in}

\usepackage{fancyhdr}
\pagestyle{fancy}
\fancyhf{}
\lhead{\CourseCode\ --- \CourseSection, \Semester}
\rhead{Homework \HomeworkNumber}
\cfoot{\thepage}
\renewcommand{\headrulewidth}{0.4pt}
\setlength{\headheight}{15pt}
\addtolength{\topmargin}{-2pt}

% =========================
% 4) Links (subtle)
% =========================
\usepackage[hidelinks]{hyperref}
\setcounter{tocdepth}{1}

% =========================
% 5) Figures (optional)
% =========================
\usepackage{graphicx}
\graphicspath{{images/}}

% =========================
% 6) Lightweight boxes (optional)
% =========================
\usepackage{mdframed}
\newmdenv[
  skipabove=0.6\baselineskip,
  skipbelow=0.6\baselineskip,
  linewidth=0.6pt,
  roundcorner=4pt,
  innerleftmargin=10pt,
  innerrightmargin=10pt,
  innertopmargin=8pt,
  innerbottommargin=8pt
]{boxnote}

% =========================
% 7) Probability / notation macros
% =========================
\newcommand{\R}{\mathbb{R}}
\newcommand{\N}{\mathbb{N}}
\newcommand{\Z}{\mathbb{Z}}

\newcommand{\OmegaSpace}{\Omega}
\newcommand{\FField}{\mathcal{F}}
% Probability measure in case it is relevant or want to distinguish it from P(A)
%\newcommand{\PMeasure}{\mathbb{P}} 
\newcommand{\E}{\mathbb{E}}

\newcommand{\given}{\,\middle|\,}
\newcommand{\ind}{\mathbf{1}}

% Common “defined as” symbol
\newcommand{\defeq}{\vcentcolon=}

% =========================
% 8) Problem command (TOC + PDF bookmarks)
% Usage:
%   \problem{<label>}
%   \problem[Short title]{<label>}
% Example:
%   \problem{1.4}
%   \problem[Sample space]{1.4}
% =========================
\makeatletter
\newcommand{\problem}[2][]{%
  \def\prob@title{Problem #2\if\relax\detokenize{#1}\relax\else\ (\textit{#1})\fi}%
  \phantomsection%
  \pdfbookmark[1]{\prob@title}{prob#2}%
  \addcontentsline{toc}{section}{\prob@title}%
  \section*{\prob@title}%
  \vspace{-0.25em}\hrule\vspace{0.8em}%
}
\makeatother

% =========================
% 9) Solution environment
% =========================
\newenvironment{solution}{%
  \noindent\textbf{Solution.}\ \ignorespaces
}{\par}

% Optional scratch / planning block
\newenvironment{scratch}{%
  \begin{boxnote}\textbf{Scratch / Setup.}\ \small
}{\end{boxnote}}

% Final answer command: expects math only (no $$ or \[ \])
\newcommand{\finalanswer}[1]{%
  \par\medskip
  \noindent\textbf{Final:}\ \fbox{$#1$}\par
} % AGAIN: MATH ONLY INSIDE THE BOX; if word answer, use \text{} inside or use \fbox{TEXT} or \framebox

% =========================
% 10) Title block
% =========================
\title{Homework \HomeworkNumber}
\author{\YourName \\
\CourseCode\ --- \CourseTitle \\
\CourseSection \\
Instructor: \Instructor}
\date{Due: \DueDate}

\begin{document}
\maketitle

% \section*{Collaboration / Resources}
% (Optional) Write “I worked alone” or list collaborators/resources.

\tableofcontents
\bigskip
\hrule
\bigskip
\newpage

% ============================================================
% Problems 
% ============================================================

\problem{2.20}
A fair die is rolled repeatedly. Use precise notation of probabilities of events and random variables for the solutions to the questions below.
\begin{enumerate}
\item Write down a precise sum expression for the probability that the first five rolls give a three at most two times.
\item Calculate the probability that the first three does not appear before the fifth roll.
\item Calculate the probability that the first three appears before the twentieth roll but not before the fifth roll.
\end{enumerate}

\begin{solution}

\begin{scratch}
Let $S_5$ be the number of threes among the first five rolls. Let $N$ be the
roll on which the first three appears. Then $S_5\sim\text{Bin}(5,1/6)$ and
$N\sim\text{Geom}(1/6)$ with support $\{1,2,\ldots\}$.
\end{scratch}

\begin{enumerate}
\item
\[
P(S_5\le 2)=\sum_{k=0}^{2}\binom{5}{k}\left(\frac{1}{6}\right)^k\left(\frac{5}{6}\right)^{5-k}.
\]
\item
\[
P(N>4)=P(\text{no three in the first four rolls})=\left(\frac{5}{6}\right)^4.
\]
\item
\[
P(5\le N\le 20)=P(\text{no three in the first four rolls})\cdot P(\text{at least one three in rolls }5\text{--}20)
\]
\[
\qquad=\left(\frac{5}{6}\right)^4\!\left(1-\left(\frac{5}{6}\right)^{16}\right)
\;=\;\left(\frac{5}{6}\right)^4-\left(\frac{5}{6}\right)^{20}.
\]
\end{enumerate}

\finalanswer{\text{(a)}\ \sum_{k=0}^{2}\binom{5}{k}\left(\frac{1}{6}\right)^k\left(\frac{5}{6}\right)^{5-k},\;
\text{(b)}\ \left(\frac{5}{6}\right)^4,\;
\text{(c)}\ \left(\frac{5}{6}\right)^4-\left(\frac{5}{6}\right)^{20}}

\end{solution}

\newpage

\problem{2.58}
Suppose that a person's birthday is a uniformly random choice from the 365 days of a year (leap years are ignored), and one person's birthday is independent of the birthdays of other people. Alex, Betty and Conlin are comparing birthdays. Define these three events:
\[
A=\{\text{Alex and Betty have the same birthday}\},
\]
\[
B=\{\text{Betty and Conlin have the same birthday}\},
\]
\[
C=\{\text{Conlin and Alex have the same birthday}\}.
\]
\begin{enumerate}
\item Are events $A$, $B$ and $C$ pairwise independent? (See the definition of pairwise independence on page 54 and Example 2.25 for illustration.)
\item Are events $A$, $B$ and $C$ independent?
\end{enumerate}

\begin{solution}

\begin{scratch}
Let birthdays be independent and uniform on $\{1,\dots,365\}$. Then
$P(A)=P(B)=P(C)=1/365$. Also $A\cap B$ means all three share a birthday.
\end{scratch}

\begin{enumerate}
\item
\[
P(A\cap B)=P(\text{all three match})=\frac{1}{365^2}
 = P(A)P(B).
\]
Similarly $P(A\cap C)=P(B\cap C)=1/365^2$, so $A,B,C$ are pairwise independent.
\item
\[
P(A\cap B\cap C)=\frac{1}{365^2}\neq \left(\frac{1}{365}\right)^3
 = P(A)P(B)P(C),
\]
so they are not mutually independent.
\end{enumerate}

\finalanswer{\text{Pairwise independent, but not mutually independent}}

\end{solution}

\newpage

\problem{2.61}
Suppose an urn has 3 green balls and 4 red balls.
\begin{enumerate}
\item Nine draws are made with replacement. Let $X$ be the number of times a green ball appeared. Identify by name the probability distribution of $X$. Find the probabilities $P(X\ge 1)$ and $P(X\le 5)$.
\item Draws with replacement are made until the first green ball appears. Let $N$ be the number of draws that were needed. Identify by name the probability distribution of $N$. Find the probability $P(N\le 9)$.
\item Compare $P(X\ge 1)$ and $P(N\le 9)$. Is there a reason these should be the same?
\end{enumerate}

\begin{solution}

\begin{scratch}
Let $p=3/7$ (green) and $q=4/7$ (red).
For nine draws with replacement, $X\sim\mathrm{Bin}(9,p)$.
For the waiting time until the first green, $N\sim\mathrm{Geom}(p)$.
\end{scratch}

\begin{enumerate}
\item $X\sim\mathrm{Bin}(9,3/7)$.
\[
P(X\ge 1)=1-P(X=0)=1-\left(\frac{4}{7}\right)^9,
\]
\[
P(X\le 5)=\sum_{k=0}^{5}\binom{9}{k}\left(\frac{3}{7}\right)^k\left(\frac{4}{7}\right)^{9-k}.
\]
\item $N\sim\mathrm{Geom}(3/7)$ with support $\{1,2,\ldots\}$.
\[
P(N\le 9)=1-P(N>9)=1-\left(\frac{4}{7}\right)^9.
\]
\item The events $\{X\ge 1\}$ and $\{N\le 9\}$ both mean
``at least one green in the first 9 draws,'' hence they are equal.
\end{enumerate}

\finalanswer{\text{(a)}\ X\sim\mathrm{Bin}(9,3/7),\ P(X\ge1)=1-(4/7)^9,\ P(X\le5)=\sum_{k=0}^{5}\binom{9}{k}(3/7)^k(4/7)^{9-k}}
\finalanswer{\text{(b)}\ N\sim\mathrm{Geom}(3/7),\ P(N\le9)=1-(4/7)^9}

\end{solution}

\newpage

\problem{2.70}
Flip a coin three times. Assume the probability of tails is $p$ and that successive flips are independent. Let $A$ be the event that we have exactly one tails among the first two coin flips and $B$ the event that we have exactly one tails among the last two coin flips. For which values of $p$ are events $A$ and $B$ independent? (This generalizes Example 2.18.)

\begin{solution}

\begin{scratch}
Let $p=P(T)$ and $q=1-p$. Then
\[
P(A)=P(\text{exactly one tail in flips 1--2})=2pq,\qquad
P(B)=2pq.
\]
Compute $P(A\cap B)$ by enumerating 3-flip sequences.
\end{scratch}

The only sequences satisfying both events are $HTH$ and $THT$, so
\[
P(A\cap B)=p q^2+p^2 q=pq.
\]
Independence requires $pq=(2pq)^2=4p^2q^2$. Thus either $pq=0$
(i.e., $p=0$ or $p=1$) or, for $0<p<1$,
\[
1=4pq \quad\Longrightarrow\quad p(1-p)=\frac14 \quad\Longrightarrow\quad p=\frac12.
\]
\finalanswer{p\in\{0,\tfrac12,1\}\ (\text{nontrivial case }p=\tfrac12)}

\end{solution}

\newpage

\problem{2.71 (a)-(c)}
As in Example 2.38, assume that 90\% of the coins in circulation are fair, and the remaining 10\% are biased coins that give tails with probability $3/5$. I hold a randomly chosen coin and begin to flip it.
\begin{enumerate}
\item After one flip that results in tails, what is the probability that the coin I hold is a biased coin? After two flips that both give tails? After $n$ flips that all come out tails?
\item After how many straight tails can we say that with 90\% probability the coin I hold is biased?
\item After $n$ straight tails, what is the probability that the next flip is also tails?
\item Suppose we have flipped a very large number of times (think number of flips $n$ tending to infinity), and each time gotten tails. What are the chances that the next flip again yields tails?
\end{enumerate}

\begin{solution}

\begin{scratch}
Let $B$ be the event the coin is biased and $F$ the event it is fair.
$P(B)=0.1$, $P(F)=0.9$, $P(T\mid B)=3/5$, $P(T\mid F)=1/2$.
For $n$ straight tails, use Bayes' rule.
\end{scratch}

\begin{enumerate}
\item For $n\ge 1$,
\[
P(B\mid T^n)=\frac{0.1(3/5)^n}{0.1(3/5)^n+0.9(1/2)^n}
 =\frac{1}{1+9(5/6)^n}.
\]
In particular,
\[
P(B\mid T)=\frac{2}{17},\qquad
P(B\mid T^2)=\frac{4}{29}.
\]
\item We need $P(B\mid T^n)\ge 0.9$, i.e.
\[
\frac{1}{1+9(5/6)^n}\ge 0.9
\;\Longleftrightarrow\;
\left(\frac{5}{6}\right)^n\le \frac{1}{81}.
\]
The smallest integer satisfying this is $n=25$.
\item The predictive probability is
\[
P(T_{n+1}\mid T^n)=\frac12\bigl(1-P(B\mid T^n)\bigr)+\frac35 P(B\mid T^n)
=\frac12+\frac{1}{10}\cdot\frac{1}{1+9(5/6)^n}.
\]
\end{enumerate}

\textit{Part (d) not assigned for HW4.}

\finalanswer{\text{(a)}\ P(B\mid T^n)=\frac{1}{1+9(5/6)^n}}
\finalanswer{\text{(b)}\ n=25}
\finalanswer{\text{(c)}\ P(T_{n+1}\mid T^n)=\frac12+\frac{1}{10}\cdot\frac{1}{1+9(5/6)^n}}

\end{solution}

\newpage

\problem{3.2}
Suppose the random variable $X$ has possible values $\{1, 2, 3, 4, 5, 6\}$ and probability mass function of the form $p(k)=ck$.
\begin{enumerate}
\item Find $c$.
\item Find the probability that $X$ is odd.
\end{enumerate}

\begin{solution}

\begin{scratch}
Normalize: $\sum_{k=1}^{6} ck = 1$.
\end{scratch}

\begin{enumerate}
\item
\[
c=\frac{1}{1+2+3+4+5+6}=\frac{1}{21}.
\]
\item
\[
P(X\ \text{odd})=c(1+3+5)=\frac{9}{21}=\frac{3}{7}.
\]
\end{enumerate}

\finalanswer{c=\frac{1}{21},\quad P(X\ \text{odd})=\frac{3}{7}}

\end{solution}

\newpage

\problem{3.4}
Let $X\sim \mathrm{Unif}[4,10]$.
\begin{enumerate}
\item Calculate $P(X<6)$.
\item Calculate $P(|X-7|>1)$.
\item For $4\le t\le 6$, calculate $P(X<t\mid X<6)$.
\end{enumerate}

\begin{solution}

\begin{scratch}
For $X\sim\mathrm{Unif}[4,10]$, $P(X\in[a,b])=(b-a)/6$.
\end{scratch}

\begin{enumerate}
\item
\[
P(X<6)=\frac{6-4}{10-4}=\frac{1}{3}.
\]
\item $|X-7|>1$ means $X<6$ or $X>8$:
\[
P(|X-7|>1)=\frac{2+2}{6}=\frac{2}{3}.
\]
\item For $4\le t\le 6$,
\[
P(X<t\mid X<6)=\frac{P(4\le X<t)}{P(4\le X<6)}=\frac{t-4}{2}.
\]
\end{enumerate}

\finalanswer{\text{(a)}\ \frac13,\ \text{(b)}\ \frac23,\ \text{(c)}\ \frac{t-4}{2}}

\end{solution}

\newpage

\problem{3.25}
In each of the following cases find all values of $b$ for which the given function is a probability density function.
\begin{enumerate}
\item
\[
f(x)=
\begin{cases}
x^2-b, & \text{if } 1\le x\le 3,\\
0, & \text{otherwise.}
\end{cases}
\]
\item
\[
h(x)=
\begin{cases}
\cos x, & \text{if } -b\le x\le b,\\
0, & \text{otherwise.}
\end{cases}
\]
\end{enumerate}

\begin{solution}

\begin{scratch}
For a p.d.f., we need nonnegativity on the support and total integral $=1$.
\end{scratch}

\begin{enumerate}
\item
\[
\int_{1}^{3}(x^2-b)\,dx=\frac{26}{3}-2b=1 \;\Longrightarrow\; b=\frac{23}{6}.
\]
But nonnegativity requires $x^2-b\ge 0$ for $x\in[1,3]$, hence $b\le 1$.
These conditions are incompatible, so no value of $b$ works.
\item
\[
\int_{-b}^{b}\cos x\,dx=2\sin b=1 \;\Longrightarrow\; \sin b=\frac12.
\]
To keep $\cos x\ge 0$ on $[-b,b]$ we need $b\le \pi/2$, so $b=\pi/6$.
\end{enumerate}

\finalanswer{\text{(a) no solution},\quad \text{(b)}\ b=\frac{\pi}{6}}

\end{solution}

\newpage

\problem{3.26 (a)}
Suppose that $X$ is a discrete random variable with possible values $\{1, 2,\ldots\}$, and probability mass function
\[
p_X(k)=\frac{c}{k(k+1)}
\]
with some constant $c>0$.
\begin{enumerate}
\item What is the value of $c$?

Hint. $\frac{1}{k(k+1)}=\frac{1}{k}-\frac{1}{k+1}$.
\end{enumerate}

\begin{solution}

\begin{scratch}
Use the telescoping sum $\frac{1}{k(k+1)}=\frac{1}{k}-\frac{1}{k+1}$.
\end{scratch}

\[
1=\sum_{k=1}^{\infty} \frac{c}{k(k+1)}
=c\sum_{k=1}^{\infty}\left(\frac{1}{k}-\frac{1}{k+1}\right)=c.
\]
\finalanswer{c=1}

\end{solution}

\newpage

\problem{3.31 (a)-(d)}
Suppose a random variable $X$ has density function
\[
f(x)=
\begin{cases}
cx^{-4}, & \text{if } x\ge 1,\\
0, & \text{else},
\end{cases}
\]
where $c$ is a constant.
\begin{enumerate}
\item What must be the value of $c$?
\item Find $P(0.5<X<1)$.
\item Find $P(0.5<X<2)$.
\item Find $P(2<X<4)$.
\end{enumerate}

\begin{solution}

\begin{scratch}
Normalize $f$ on $[1,\infty)$ and integrate for the probabilities.
\end{scratch}

\begin{enumerate}
\item
\[
1=\int_{1}^{\infty} c x^{-4}\,dx=c\cdot\frac{1}{3}\quad\Longrightarrow\quad c=3.
\]
\item Since the support is $[1,\infty)$,
\[
P(0.5<X<1)=0.
\]
\item
\[
P(0.5<X<2)=\int_{1}^{2}3x^{-4}\,dx
=3\left[\frac{-1}{3}x^{-3}\right]_{1}^{2}=\frac{7}{8}.
\]
\item
\[
P(2<X<4)=\int_{2}^{4}3x^{-4}\,dx
=3\left[\frac{-1}{3}x^{-3}\right]_{2}^{4}=\frac{7}{64}.
\]
\end{enumerate}

\finalanswer{\text{(a)}\ c=3,\ \text{(b)}\ 0,\ \text{(c)}\ \frac{7}{8},\ \text{(d)}\ \frac{7}{64}}

\end{solution}


\end{document}
